\documentclass[11pt,a4paper]{article}
\usepackage[margin=20mm]{geometry}
\usepackage{fontspec}
\setmainfont{Times New Roman}
\setmonofont{Courier New}
\usepackage[russian]{babel}
\usepackage{listings}
\setlength{\parindent}{0pt}
\setlength{\parskip}{5pt}
\lstset{basicstyle=\ttfamily\fontsize{8.2}{9.5}\selectfont,
  columns=fullflexible,keepspaces=true,showstringspaces=false,
  breaklines=true,numbers=none,aboveskip=4pt,belowskip=6pt}
\newcommand{\head}[1]{\par\medskip{\large\bfseries #1}\par\nobreak\smallskip}
\newcommand{\subhead}[1]{\par\smallskip{\bfseries #1}\par\nobreak\smallskip}
\begin{document}
\thispagestyle{empty}
\vspace*{40mm}
\begin{center}
{\Large\bfseries ОТЧЁТ}\par
по лабораторной работе №2\par
по дисциплине «Операционные системы»\par
\vspace{8mm}
{\large Процессы и неименованные каналы}\par
Вариант 4
\end{center}
\vspace{35mm}
\begin{flushright}
Выполнил: Григорьев Глеб Алексеевич\par
Группа: М8О-202БВ-25
\end{flushright}
\vfill
\begin{center}2026\end{center}
\newpage
\head{1 Цель работы}
Изучить создание дочернего процесса и обмен данными через неименованные каналы. Использовать перенаправление стандартных потоков и ожидание завершения дочернего процесса.

\head{2 Постановка задачи}
Родительская программа получает имя выходного файла и строки с числами типа float. Дочерняя программа последовательно делит первое число строки на остальные и записывает результат в файл. При нулевом делителе оба процесса должны завершиться.

\head{3 Организация программы}
Решение состоит из программ parent и child. Общие функции обмена находятся в ipc.c, их объявления и коды состояний — в ipc.h. Для сборки используется Makefile.

Родитель создаёт два канала до fork. Через pipe1 он передаёт строки в стандартный ввод дочернего процесса. Через pipe2 дочерний процесс возвращает состояние обработки. Имя выходного файла передаётся аргументом при запуске child через execl.

После каждой строки родитель ждёт ответ: K — результат записан, Z — нулевой делитель, E — ошибка. Перед началом работы дочерний процесс отправляет R, если файл успешно открыт. Такой порядок позволяет завершить оба процесса при делении на ноль без следующего ввода пользователя.

\head{4 Основные фрагменты программы}
\subhead{Создание каналов и запуск дочернего процесса — parent.c}
\lstinputlisting[firstline=45,lastline=70]{../parent.c}
\newpage
\head{4 Основные фрагменты программы (продолжение)}
\subhead{Преобразование чисел и деление — child.c}
Первое число сохраняется в result. Каждый следующий делитель проверяется на ноль до выполнения операции. strtof позволяет проверить, что число прочитано полностью и находится в допустимом диапазоне.

\lstinputlisting[firstline=21,lastline=41]{../child.c}
\subhead{Запись в канал — ipc.c}
write может передать только часть строки. Функция повторяет запись оставшихся байтов, а при EINTR повторяет прерванный вызов.

\lstinputlisting[firstline=7,lastline=23]{../ipc.c}
SIGPIPE игнорируется, поэтому обрыв канала обрабатывается как ошибка write. При закрытии канала ответов родитель получает EOF. Ошибки системных вызовов проверяются, сообщения выводятся в stderr.

Перед завершением освобождаются буферы getline, закрываются файл и дескрипторы каналов. Родитель вызывает waitpid. При обычном конце ввода запись в pipe1 закрывается, дочерний процесс получает EOF и завершает работу.

\newpage
\head{5 Результаты проверки}
Проверка выполнена 18 ноября 2026 года в контейнере Debian 12, Linux aarch64. Использовались GCC 12.2, Python 3.11 и strace 6.1. Команды сборки и тестирования:

\begin{lstlisting}
make clean
make
make test
\end{lstlisting}
Сборка с -Wall -Wextra -Wpedantic -Werror завершилась без предупреждений. Прошли все 18 тестов, включая запись в /dev/full, длинные строки, ошибочный ввод и завершение по нулевому делителю.

\subhead{Обычный запуск}
Для строк «100 2 5», «7.5 2.5» и «0 2» в файле получены 10, 3 и 0. Ниже приведены строки журнала normal.log с PID процессов:

\begin{lstlisting}
[parent PID=890] Output file: [parent PID=890] child PID=891
[child PID=891 PPID=890] Started
[parent PID=890] Numbers (Ctrl+D to finish): [child PID=891] Result: 10
[parent PID=890] Numbers (Ctrl+D to finish): [child PID=891] Result: 3
[parent PID=890] Numbers (Ctrl+D to finish): [child PID=891] Result: 0
[parent PID=890] Numbers (Ctrl+D to finish): [parent PID=890] Child exit code: 0
\end{lstlisting}
\subhead{Нулевой делитель}
На вход поданы «20 2», «10 0» и «99 3». В файле остался только результат 10. Третья команда не выполнена, оба процесса завершились с кодом 2. Фрагмент zero.log:

\begin{lstlisting}
[parent PID=897] Output file: [parent PID=897] child PID=898
[child PID=898 PPID=897] Started
[parent PID=897] Numbers (Ctrl+D to finish): [parent PID=897] Division by zero; stopping
[parent PID=897] Child exit code: 2
\end{lstlisting}
\subhead{Трассировка системных вызовов}
Запуск выполнялся через strace -f -o /logs/normal.strace ./parent. Фрагменты normal.strace показывают создание каналов, запуск дочернего процесса, запись результата и ожидание его завершения:

\begin{lstlisting}
890   pipe2([3, 4], 0)                  = 0
890   pipe2([5, 6], 0)                  = 0
890   clone(child_stack=NULL, flags=CLONE_CHILD_CLEARTID|CLONE_CHILD_SETTID|SIGCHLD, child_tidptr=0xffff80d4b050) = 891
891   dup3(3, 0, 0 <unfinished ...>
891   <... dup3 resumed>)               = 0
891   dup3(6, 1, 0 <unfinished ...>
891   <... dup3 resumed>)               = 1
891   execve("./child", ["child", "/logs/normal-results.txt"], 0xffffe0dd7bc8 /* 6 vars */) = 0
891   write(3, "10\n", 3)               = 3
891   exit_group(0)                     = ?
890   <... wait4 resumed>[{WIFEXITED(s) && WEXITSTATUS(s) == 0}], 0, NULL) = 891
890   exit_group(0)                     = ?
\end{lstlisting}
В этой Linux-среде вызовы fork, dup2 и waitpid видны как clone, dup3 и wait4. Полные журналы и результаты тестов сохранены в docs/assets/linux.

\head{6 Вывод}
Реализованы две программы, обменивающиеся строками и состояниями через неименованные каналы. Деление выполняется в дочернем процессе. Проверены обычное завершение, нулевой делитель и ошибочный ввод. Linux-тесты и трассировка подтвердили работу обмена и завершение обоих процессов.

\end{document}
